\subsection{Processus expérimental physique}

\subsubsection{Chaîne de mesure APEX : Acquisition et transmission de données inertielle}

\newpage

\subsubsection{APEX BES : Banc d’essai pour la validation de la chaîne APEX}
\label{sec:333_banc_test_apex}

La validation de la chaîne de mesure et de transmission d’APEX s’appuie sur un banc d’essai dédié, développé pour reproduire au
sol des profils de mouvement représentatifs et répétables. Le cœur de ce travail est la conception et la fabrication de deux
véhicules d’essai capables d’exécuter une trajectoire de type parabolique (au sens du profil cinématique visé), afin de
solliciter l’acquisition inertielle et la transmission radio dans un cadre maîtrisé. La capture vidéo en vue zénithale est
utilisée comme \textit{référence externe} pour comparer et interpréter les mesures issues d’APEX. Le proejt APEX BES (Banc d’Essai
SP-02) constitue ainsi un sous-système d’essai complet, permettant de valider la chaîne APEX dans des conditions proches de celles
rencontrées en vol.\\

\subsubsubsection*{Objectif de l’essai}
Le banc vise à :
\begin{itemize}
    \item \textbf{générer une cinématique reproductible} (profil de déplacement et changements de régime identifiables),
    \item \textbf{valider la chaîne APEX} (acquisition \acrshort{IMU}, enregistrement, transmission et affichage station sol),
    \item \textbf{corréler les mesures} APEX avec une trajectoire « vraie » reconstruite par suivi vidéo, afin de disposer
    d’un outil d’analyse pour qualifier la cohérence globale (ordre de grandeur, synchronisation, détectabilité).
\end{itemize}

\subsubsubsection*{Conception et réalisation des véhicules d’essai}
Deux véhicules ont été conçus et fabriqués spécifiquement pour ce banc. Ils servent de supports instrumentés pour embarquer
l’électronique APEX, tout en permettant d’imposer une trajectoire et des excitations dynamiques contrôlées. Les choix de
conception visent à garantir :
\begin{itemize}
    \item une \textbf{répétabilité} des essais (géométrie et masse connues, comportement stable),
    \item une \textbf{facilité d’intégration} des cartes (fixations, accessibilité, protection),
    \item une \textbf{cinématique exploitable} (accélérations et rotations suffisamment marquées pour émerger du bruit).
\end{itemize}
Au-delà d’un simple support, ces véhicules constituent un sous-système d’essai à part entière : leur fabrication a permis de
maîtriser les interfaces mécaniques (montage des cartes, passage câbles, rigidité locale) et de disposer d’un moyen autonome
pour reproduire des scénarios de mouvement.

% -----------------------------------------------------------------
% FIGURES : à remplacer par les figures du rapport APEX BES
% -----------------------------------------------------------------
\begin{figure}[H]
    \centering
    % \includegraphics[width=0.92\textwidth]{figures/apex_bes/fig_01.png}
    \caption{%
    % CAPTION À COPIER-COLLER du rapport APEX BES (Figure 1)
    }
    \label{fig:apex_bes_fig01}
\end{figure}

\begin{figure}[H]
    \centering
    % \includegraphics[width=0.92\textwidth]{figures/apex_bes/fig_02.png}
    \caption{%
    % CAPTION À COPIER-COLLER du rapport APEX BES (Figure 2)
    }
    \label{fig:apex_bes_fig02}
\end{figure}

\subsubsubsection*{Protocole d’essai}
Les essais sont conduits sur une fenêtre temporelle commune en enregistrant simultanément : (i) les données APEX (\acrshort{IMU},
télémesure) et (ii) une vidéo zénithale permettant de reconstruire la trajectoire. Le déroulement est le suivant :
\begin{enumerate}
    \item mise en place des véhicules et vérification des intégrations APEX,
    \item calibration de la scène (repère sol et échelle pixel $\rightarrow$ mètre),
    \item exécution de la trajectoire par les véhicules (profil parabolique visé) et enregistrement,
    \item traitement des données vidéo (trajectoire puis dérivation) et extraction des courbes APEX,
    \item alignement temporel et comparaison des grandeurs.
\end{enumerate}

\subsubsubsection*{Trajectoire de référence par capture vidéo}
La capture vidéo (vue zénithale) fournit une référence géométrique externe. La trajectoire est obtenue en suivant les positions
$(x(t),y(t))$ des véhicules puis en appliquant la conversion d’échelle. Les vitesses et accélérations sont calculées par
dérivation numérique sur un pas $\Delta t$ correspondant à la cadence caméra :
\[
v_x(t) \simeq \frac{x(t+\Delta t)-x(t)}{\Delta t}, \qquad
a_x(t) \simeq \frac{v_x(t+\Delta t)-v_x(t)}{\Delta t}
\]
(analogues pour $y$). Un lissage léger peut être appliqué à la trajectoire avant dérivation afin de limiter l’amplification du
bruit, tout en conservant la dynamique utile.

\subsubsubsection*{Alignement temporel}
Les données APEX et la vidéo ne partagent pas la même horloge. L’alignement est réalisé à partir d’un événement commun
identifiable (démarrage de mouvement, arrêt net, ou excitation volontaire générant un pic). Cette étape conditionne la lisibilité
des comparaisons, en particulier sur les dérivées (accélérations).

\begin{figure}[H]
    \centering
    % \includegraphics[width=0.92\textwidth]{figures/apex_bes/fig_03.png}
    \caption{%
    % CAPTION À COPIER-COLLER du rapport APEX BES (Figure 3)
    }
    \label{fig:apex_bes_fig03}
\end{figure}

\newpage